\section{Threat Model}
\label{sec:threatmodel}

In this paper, we focus on malicious code poisoning attacks within the AI model supply chain, in which malicious actors exploit the collaborative nature of pre-trained model hubs to distribute compromised artifacts and compromise victim systems.

\noindent \textbf{Attackers.} The attackers are malicious AI model providers or actors who have compromised legitimate developer accounts\cite{10754708}. They have the capability to create and upload malicious datasets and binary model files to public model hubs without providing the original training source code. To embed their payloads, attackers can exploit insecure serialization formats\cite{liu2025arthideseekmaking} or abuse legitimate, hidden capabilities of framework APIs\cite{zhu2025} to bypass static security scanners. They may also employ techniques like exploiting leaked authentication tokens or hijack trusted identities, thereby deceiving the community\cite{3773117}. Their primary goal is to execute unauthorized actions, such as establishing reverse shells, exfiltrating sensitive data such as user credentials, or deploying ransomware, without attracting attention.

\noindent \textbf{Victim Users.} The victims are AI developers and researchers who seek to download and reuse pre-trained models or datasets for their applications\cite{stalnaker2025empiricalanalysismachinelearning}. They generally trust content hosted on well-known model hubs and popular contributors, often prioritizing convenience and efficiency over rigorous security inspections\cite{siddiq2026empiricalstudyremotecode}. Victims typically load and execute these models using standard libraries and standard user-level permissions, rather than sandboxed environments. During the standard model loading or inference processes, they unknowingly trigger the embedded malicious code or behaviors\cite{models_are_codes}.

\noindent \textbf{Defenders:} The defenders deploy \toolname{} to protect their systems against malicious models. \toolname{} operates in a controlled sandbox environment and is designed for two deployment scenarios. In the primary scenario, \toolname{} is integrated into a pre-deployment model vetting pipeline: before a model is released or integrated into any production system, it undergoes security vetting where \toolname{} monitors the model's behavior during loading and inference to detect embedded malicious payloads. To maximize coverage, the model should be exercised through all anticipated usage scenarios—including loading with expected input shapes and data types, inference across supported batch sizes, and edge-case inputs that stress runtime boundaries—ensuring that conditionally triggered malicious logic is activated and exposed during testing\cite{11029791, zhang2025fixexploitenablingcomprehensive}. In the secondary scenario, for environments where inference latency is not a critical constraint, \toolname{} can be deployed directly at model serving time to continuously monitor model behavior and block suspicious execution in real time.